\documentclass[10pt]{article}

\usepackage[french]{babel}
\usepackage[T1]{fontenc}
\usepackage[utf8]{inputenc}
\usepackage{amsmath,amssymb}

\begin{document}

\section*{Orthogonalité des fonctions propres}

On considère le problème spectral

\[
X_n''(u)+\lambda_n^2X_n(u)=0
\text{\quad sur l'intervalle \quad} u\in\left[-\frac12,\frac12\right],
\]

avec les conditions de Robin homogènes:

\[
X_n'\!\left(-\frac12\right)
=
Bi\,X_n\!\left(-\frac12\right)
\text{\quad et \quad}
X_n'\!\left(\frac12\right)
=
-Bi\,X_n\!\left(\frac12\right).
\]

Les valeurs propres $\lambda_n$ sont déterminées par:

\[
\tan\left(\frac{\lambda_n}{2}\right)
=
\frac{Bi}{\lambda_n}.
\]

Nous montrons que les fonctions propres associées à deux valeurs propres distinctes sont orthogonales.

\subsection*{Étape 1 : Équations différentielles}

Les fonctions propres $X_n$ et $X_p$ vérifient

\[
X_n''+\lambda_n^2X_n=0,
\text{\quad et \quad}
X_p''+\lambda_p^2X_p=0.
\]

On multiplie la première équation par $X_p$ et la seconde par $X_n$, puis on soustrait :

\[
X_pX_n''
-
X_nX_p''
+
(\lambda_n^2-\lambda_p^2)X_nX_p
=
0.
\]

Or,

\[
X_pX_n''
-
X_nX_p''
=
\frac{d}{du}
\left(
X_pX_n'
-
X_nX_p'
\right),
\]

d'où

\[
\frac{d}{du}
\left(
X_pX_n'
-
X_nX_p'
\right)
+
(\lambda_n^2-\lambda_p^2)X_nX_p
=
0.
\]

\subsection*{Étape 2 : Intégration}

On intègre cette relation sur l'intervalle
$\left[-\frac12,\frac12\right]$ :

\[
(\lambda_n^2-\lambda_p^2)
\int_{-1/2}^{1/2}
X_nX_p\,du
=
-
\left[
X_pX_n'
-
X_nX_p'
\right]_{-1/2}^{1/2}.
\]

Il reste donc à évaluer le terme de bord.

\subsection*{Étape 3 : Conditions de Robin}

Au bord droit,

\[
X_n'
=
-Bi\,X_n,
\qquad
X_p'
=
-Bi\,X_p.
\]

Ainsi,

\[
X_pX_n'
-
X_nX_p'
=
-Bi\,X_pX_n
+
Bi\,X_nX_p
=
0.
\]

Au bord gauche,

\[
X_n'
=
Bi\,X_n,
\qquad
X_p'
=
Bi\,X_p,
\]

d'où

\[
X_pX_n'
-
X_nX_p'
=
Bi\,X_pX_n
-
Bi\,X_nX_p
=
0.
\]

Les deux contributions au terme de bord sont donc nulles.

On obtient alors

\[
(\lambda_n^2-\lambda_p^2)
\int_{-1/2}^{1/2}
X_nX_p\,du
=
0.
\]

Comme les valeurs propres sont distinctes,

\[
\lambda_n\neq\lambda_p,
\]

on conclut que

\[
\boxed{
\int_{-1/2}^{1/2}
X_n(u)X_p(u)\,du
=
0,
\qquad
n\neq p.
}
\]

\section*{Conclusion}

Les fonctions propres associées aux conditions de Robin homogènes forment une famille orthogonale pour le produit scalaire usuel

\[
\langle X_n,X_p\rangle
=
\int_{-1/2}^{1/2}
X_n(u)X_p(u)\,du.
\]

Cette propriété est une conséquence directe du caractère auto-adjoint de l'opérateur de Sturm--Liouville

\[
L[X]=-X''.
\]

\end{document}